\documentclass{article}
\usepackage[margin=3cm]{geometry}
\usepackage{amsmath,amssymb,mathtools}
\usepackage{longtable,array}
\usepackage{hyperref}
\hypersetup{hidelinks}
\pagestyle{plain}

% Objects and fields.
\newcommand{\genesis}{B_{\mathrm{gen}}}
\newcommand{\slot}{\operatorname{slot}}
\newcommand{\depth}{\operatorname{depth}}
\newcommand{\rank}{\operatorname{rank}}
\newcommand{\T}{\mathcal T}
\newcommand{\AC}{\mathcal A}
\newcommand{\Fr}{\Phi}
\newcommand{\tree}{\mathsf{tree}}
\newcommand{\fcroot}{\mathsf{fc\_root}}
\newcommand{\gfPool}{\mathsf{gf\_votes}}
\newcommand{\fname}[1]{\text{\textsc{#1}}}
\newcommand{\latestConfirmed}{\mathsf{latest\_confirmed}}
\newcommand{\timestamp}{\mathsf{timestamp}}
\newcommand{\savedConfirmation}{\mathsf{saved\_confirmation}}
\newcommand{\locks}{\mathsf{locks}}
\newcommand{\restartTime}{\mathsf{restart\_time}}
\newcommand{\acks}{\mathsf{acks}}
\newcommand{\acked}{\mathsf{acked}}

\usepackage{xurl}
\newcommand{\nm}[1]{\nolinkurl{#1}}
\newcommand{\Lc}{\textbf{Lean change}}
\newcommand{\Sc}{\textbf{spec change}}
\newcommand{\Ld}{\textbf{lead}}
\newcommand{\Pm}{\textbf{premise}}
\newcommand{\Fm}{\textbf{form}}
\newcommand{\UC}{\textbf{UNCHECKED}}

\title{PC-SA specification: relation to the stage-2 Lean model}
\author{}
\date{}
\begin{document}
\maketitle

This note is a companion to \nm{pcsa-spec.tex} on the branch
\nm{pcsa-rebase}, after the stage-2 port (ruling R-spec-a).  It maps the
names of the specification to the stage-2 Lean model, it lists the
differences with their status, and it maps the claims.  Its section
numbers refer to this note.  ``Spec~3.7'' is Section~3.7 of the
specification, and ``S32'' is a row of its decision ledger.

\emph{The Lean of record.}  The repository is
\nm{~/ethereum/lean/available-chain}, branch \nm{pc-int16} (tip
\nm{45926c3b2}; \nm{pc-int15} is \nm{ac4745e99}).  It was read on
2026-10-19 from a read-only worktree.  No Lean file was changed and
nothing was built.  The release model is stage~2:
\begin{itemize}
\item M: \nm{pcsa/AvailableChainPCSA/Model/Stage2/}; a citation
  ``Gate:34'' is line~34 of \nm{Gate.lean} in that folder.
\item P: \nm{pcsa/AvailableChainPCSARun/Model/Stage2/Protocol.lean}, the
  run protocol \nm{pcProtocol} (P:1241) on the engine
  \nm{V2.ProtocolSpecD}; ``P:996'' is a line of that file.
\item The run driver: \nm{Driver.lean} and \nm{Macro.lean} in the folder
  of P (``Driver:383'', ``Macro:171'').
\item The claims:
  \nm{pcsa/AvailableChainPCSAClaims/Statements/Stage2/Claims.lean}.
\end{itemize}
Lean names have the namespace \nm{AvailableChain.PCSA.Stage2}.  The model
follows \nm{drafts/pcsa-n2-r4-stage2.md} (S2); its docstrings cite S2
lines.  The old clone \nm{~/lean/available-chain} and its stage-1 model
are legacy (Section~\ref{sec:legacy}).

\emph{Status words.}  A \Fm{} difference changes no output and no message
of an honest node, under the stated condition.  A behaviour difference
changes a rule, a wire object, a read, or the scope of a claim.  For each
one the note says what is necessary: a \Lc{} (not made), a \Sc{}, a
\Pm{} of the claims, or a decision of the \Ld.  A statement that this note
did not check against the Lean code is marked \UC.

\section{Names}\label{sec:names}

The rows follow the figures of the specification.

{\footnotesize
\setlength{\tabcolsep}{3pt}
\begin{longtable}{@{}>{\raggedright\arraybackslash}p{4.6cm}>{\raggedright\arraybackslash}p{8.4cm}@{}}
\hline
Specification & Lean (M unless P)\\
\hline
\endhead
\multicolumn{2}{@{}l}{\emph{Spec~3.1, schedule and parameters}}\\
$t_s,v_s,e_s,z_s,q_s,d_s$, $a_r$ & \nm{proposalTime}, \nm{voteTime}, \nm{earlyTime}, \nm{freezeTime}, \nm{releaseTime}, \nm{gateTime}, \nm{actionTime} (Params:42--49)\\
$k$, $L$, $\mu$ & \nm{ℓ}, \nm{L}, \nm{μ} $=$ \nm{pcQuorum n} (Params:33--35, PCKey:37)\\
\multicolumn{2}{@{}l}{\emph{Spec~3.2, state (Table~1)}}\\
$\gfPool$, vote stamps & P: \nm{core.heads}, \nm{core.headReceipts}, \nm{core.headResolutions}\\
block stamps $\timestamp[B]$ & P: \nm{core.bodyStamps}\\
$\acks$ & the ack pool of \nm{onAck} (Pools:43)\\
$\savedConfirmation$, $\mathsf{gate\_ok}$, $\mathsf{ft}$ & \nm{rawResults}, the flags \nm{(w, guard, ft)} of \nm{runGate}\\
$\mathsf{ack\_ok}$ & the third output of \nm{prepareClock}; P: \nm{core.ackEnabled}\\
$\locks$, $\latestConfirmed$ & \nm{locks}, \nm{latestConfirmed} (State:29--45)\\
$M$, a role pair & \nm{maximum}, \nm{RolePair} (State:10)\\
$\mathsf{early}$, $\mathsf{recent}$, $\mathsf{unknown\_recent}$, $\mathsf{interval}$ & \nm{early}, \nm{recent} (\nm{Recent}, State:21), \nm{unknownRecent}, \nm{interval}\\
$\mathsf{freeze}[2]$, $\mathsf{carried}[2]$ & \nm{freeze0}, \nm{freeze1}, \nm{carrier0}, \nm{carrier1} (\nm{LabelledRole}, State:15)\\
$\mathsf{fg\_candidate}$ & \nm{sourceCandidate}\\
$\mathsf{sent\_keys}$ & P: \nm{core.firstKeys}\\
$\acked$ & none (S33)\\
$\restartTime$ & \nm{restartTime}\\
\multicolumn{2}{@{}l}{\emph{Spec~3.3, messages and handlers (Fig.~3)}}\\
vote, ack, PC & \nm{HeadWire}, \nm{AckWire}, \nm{PC} (Signing:41, PCKey:41, PCKey:56)\\
key, rank, valid & \nm{PCKey}, \nm{PC.rank}, \nm{rankBelow}, \nm{preferPC}, \nm{PC.valid} (PCKey:10--84)\\
payload, digest input & \nm{NamedProposal}, \nm{DigestInput} (ProposalContent:26, 115)\\
\fname{ac\_valid\_payload} & P: \nm{bodyValid} (P:268), \nm{namedContentOk} (Objects:9)\\
\fname{ac\_on\_tick} & \nm{prepareClock} (Clock:25); P: \nm{tick} (P:996)\\
\fname{ac\_duties} & P: \nm{act} duties 1 and 3 (P:1126)\\
\fname{on\_goldfish\_vote} & \nm{onHead} (Pools:13)\\
\fname{on\_ack} & \nm{onAck} (Pools:43)\\
\fname{ac\_on\_block} & \nm{ac\_on\_block} (BlockCallback:35), after \nm{pcLeaf} (Leaf:17)\\
\multicolumn{2}{@{}l}{\emph{Spec~3.4, fork choice (Fig.~4)}}\\
\fname{resolution\_time}, \fname{vote\_view} & P: \nm{resolveHeads}, \nm{batchBefore}, \nm{batch}, \nm{releaseBatch}, \nm{outputBatch}\\
\fname{equivocators}, \fname{supporters}, \fname{opponents}, \fname{goldfish\_score} & \nm{equivocatingSeats}, \nm{supportCount}, \nm{oppositionCount}, \nm{dutyScore} (Walks:11--38)\\
\fname{latest\_view} & \nm{outputBatchLabel} (Outputs:30) with P: \nm{outputBatch}\\
\fname{ghost}, \fname{get\_head} & \nm{dutyWalk} (Walks:69)\\
\fname{viable\_tree} & \nm{activeDomain}, \nm{viableDomain} (Frame:11--15)\\
\fname{status}, \fname{authority} & \nm{targetStatus}, \nm{roleAuthority} (Authority:16--25)\\
\fname{anc} & \nm{latestAncestor}, \nm{authorityAncestor} (Authority:31--36)\\
\fname{get\_anchor} & \nm{framePrefix}, \nm{framedBase}, \nm{liveLockRoot} (Frame:25--44)\\
\fname{ac\_head} & \nm{outputHead} (Outputs:67)\\
\multicolumn{2}{@{}l}{\emph{Spec~3.5, locks (Fig.~5)}}\\
\fname{effective} & \nm{effectiveLock} (Frame:38), \nm{lockCompatible} (History:34)\\
\fname{unlock} & \nm{releaseLocks} (Frame:51) with P: \nm{releaseBatch}\\
\fname{write\_lock} & the lock fold of \nm{runGate} (Gate:10)\\
\multicolumn{2}{@{}l}{\emph{Spec~3.6, confirmation (Fig.~6)}}\\
\fname{get\_confirmation} & \nm{rawConfirmation} (Walks:43)\\
\fname{update\_confirmation} & \nm{runGate} (Gate:10), \nm{saveRawResult}, \nm{slowWindow}, \nm{slowWriteAllowed}, \nm{rawLockGuard} (Walks:51--64)\\
\fname{history\_ok}, \fname{known\_compatible} & \nm{boundedHistoryCompatible}, \nm{historyCompatible} (History:29, 69)\\
\fname{outputs}, \fname{ac\_strong}, \fname{ac\_confirmed} & \nm{pcOutputs}, \nm{normalizeConfirmed}, \nm{readOutputs} (Outputs:17--79)\\
\multicolumn{2}{@{}l}{\emph{Spec~3.7, pre-confirmations (Fig.~7)}}\\
\fname{role\_receipt} & \nm{roleReceipt} (Roles:14)\\
\fname{on\_pc} & \nm{on\_pc} (PCProcess:14), \nm{historyReceipt} (History:75), \nm{carrierReceipt} (Roles:47)\\
\fname{add\_recent}, \fname{recent\_compatible} & \nm{addRecent}, \nm{recentCompatible} (History:44--61)\\
\fname{try\_assemble} & \nm{assemblePC} (Pools:29), called by \nm{onAck}\\
\fname{target\_accepted} & \nm{notifyAccepted}, \nm{acceptTargets}, \nm{roleAccepted} (BlockCallback:10, Roles:20--24)\\
\fname{early\_cut}, \fname{save\_freeze}, \fname{open\_carrier} & \nm{earlyCut} (History:85), \nm{saveFreeze}, \nm{openCarrier} (Roles:33--42)\\
\fname{view\_source}, \fname{recent\_source} & \nm{viewSource}, \nm{recentSource} (Roles:56--63)\\
\fname{ack\_test} & \nm{ackEligible} (Gate:34); \nm{recentAllowsAck} (History:39)\\
\multicolumn{2}{@{}l}{\emph{Spec~3.8, duties (Fig.~8)}}\\
\fname{voter\_block\_tree}, \fname{merged\_vote\_view} & P: \nm{headDomain}, \nm{headBatch} (P:910--924)\\
\fname{goldfish\_vote} & P: \nm{headVote} (P:949)\\
\fname{ac\_proposal} & P: \nm{act} duty 0, \nm{paired} (P:966)\\
\fname{prepare\_fg\_source}, \fname{ac\_to\_fg} & \nm{prepareSource}, \nm{ac\_to\_fg} (Source:11--27); P: \nm{fgInputs} (P:979)\\
\multicolumn{2}{@{}l}{\emph{Spec~3.9, joining and restart (Fig.~9)}}\\
\fname{joined} & \nm{clockJoined} (Clock:10)\\
\fname{ac\_restart} & \nm{restorePC}, \nm{reconstructMissedFreezes}, \nm{ac\_restart} (Restart:12--75); \nm{scanCumulativeCarriers}, \nm{normalizeRestartCompanions} (Replay:31--40); P: \nm{Durable}, \nm{durable}, \nm{restore} (P:1018--1045); Driver: \nm{reconstruct}, \nm{replays}, \nm{replay}, \nm{rebuild}; Macro: \nm{restartWork}, \nm{macroRestore}\\
\multicolumn{2}{@{}l}{\emph{Spec~4, PC-SA node (Fig.~10)}}\\
\fname{on\_tick} & P: \nm{tick}, \nm{due}, \nm{act}, \nm{recordOwn} (P:996--1188); \nm{selectDuties} (Clock:58)\\
\fname{on\_block}, dependency service & P: \nm{processOrdinary}, \nm{processReceipt}, \nm{pendingTokens}, \nm{refusalEntries}, \nm{targetPath} (P:363--790); \nm{pcLeaf}\\
\fname{pc\_frame} & P: \nm{readTarget}, \nm{refusedRoot}, \nm{tree}, \nm{controls}, \nm{floor}, \nm{frameRoot}, \nm{allowed} (P:799--851)\\
relay & the relay actions of \nm{onHead}, \nm{onAck}, \nm{on\_pc}, \nm{ac\_on\_block}; P: \nm{consumeEffects} (P:566)\\
\hline
\end{longtable}
}

\section{Differences}\label{sec:diff}

The specification follows S2 and the Lean.  Where they differ, the
ledger of the specification says which form it takes.  The table lists
every difference that this note found.  Items (a), (c)--(e), (g), (k)
and (m) are \Fm.  By ruling R-spec-b the specification follows the Lean
code unless the Lean behaviour is a defect; items (f), (h)--(j), (l), (n)
and (p) are now equal to the Lean, and each paragraph says what differs
from S2 or from Goldilocks.

\begin{table}[!ht]
\centering
\footnotesize
\caption{The differences between the specification and the stage-2 Lean.}\label{tab:diff}
\begin{tabular}{@{}l>{\raggedright\arraybackslash}p{4.0cm}>{\raggedright\arraybackslash}p{4.2cm}>{\raggedright\arraybackslash}p{2.4cm}@{}}
\hline
Item & Specification & Lean & Status\\
\hline
(a) S13 & liveness test $d>t$ & expiry filter at each tick & form\\
(b) S11, S44 & no application value; parent root in the digest & \nm{app}; full parent in \nm{DigestInput} & content abstraction\\
(c) S24 & $\Fr.\T$, $\Fr.\mathsf{refused}$ in a PC frame & \nm{read}, \nm{refused} arguments & form\\
(d) V3 & carried PC with the votes & \nm{onHead} on \nm{raw ++ support} & form\\
(e) S21, S33 & $\AC.\acked$ and A3 & no record; one duty tick & form (one tick per time)\\
(f) S35 & output head with the current duty slot & the same; S2 has $s^*+1$ & ruled; spec $=$ Lean $\neq$ S2\\
(g) S36 & 13 PC records & \nm{State.pcSlots} & form\\
(h) S17, S45 & select; own objects at delivery & \nm{selectDuties}; \nm{recordOwn} sends only & ruled; spec $=$ Lean\\
(i) S32 & A2 reads the gate batch of slot $w-1$ & \nm{batch (s - 2)} & ruled; spec $=$ Lean; docstring defect\\
(j) S34 & no slot-0 duty & \nm{0 < s} in \nm{clockJoined} & ruled; spec $=$ Lean\\
(k) S40 & restart order of S2 \S5.1 & the same, in the run driver & form (UNCHECKED in part)\\
(l) S39 & LATEST-ANC & LATEST-ANC; D21c has R-OPP3 & ruled; spec $=$ Lean; research OPEN\\
(m) S2 \S6.1 & dependency service in the node & driver \nm{drainState}; P durable fields & form; shared item (e)\\
(n) S41 & no slot test in resolution; views keep it & the same & ruled; spec $=$ Lean\\
(o) S42 & one population, $n=N$ & PC seats $n$, FG population $N$ & specialization\\
(p) S43 & witnesses by seat index & the same & ruled; spec $=$ Lean\\
\hline
\end{tabular}
\end{table}

\emph{(a) Lock expiry.}  \nm{prepareClock} first removes each lock with
$t\geq d$ (Clock:29).  The specification keeps a dead lock and tests
$d>t$ at each use.  A dead lock changes no decision in either form, and
\nm{releaseLocks} and \nm{effectiveLock} test the deadline.  \Fm.

\emph{(b) Payload content.}  \nm{NamedProposal} has \nm{app}
(ProposalContent:48), and \nm{act} duty~0 puts it in an honest proposal;
\nm{DigestInput} (ProposalContent:115--126) has \nm{app} and the full
parent, where the specification has no application value (S11) and the
parent root.  No AC rule and no claim reads \nm{app}.  So the honest
blocks differ in their bytes and roots, and the difference is not \Fm{}
under the definition above.  The correspondence is a content abstraction:
a Lean block maps to a block of the specification by dropping \nm{app}
and by the root of its parent (S44).

\emph{(c) Frame.}  The Lean passes \nm{read} (the cumulative accepted
bodies, \nm{historicalTarget}) and \nm{refused} to each function.  The
specification gives the same reads as two fields of the PC frame.  The
effective-lock test of the Lean is $\Fr.F\preceq B$ and $B$ held; the
specification writes $B\in\Fr.\tree$, which is the same set, because
$\Fr.\tree$ is the set of accepted blocks that extend $\Sigma.F$.  \Fm.

\emph{(d) Carried votes.}  \nm{ac\_on\_block} admits \nm{raw ++ support}
(BlockCallback:47--50); the support list is a subset of the raw list
(\nm{bodyValid}), so the admitted set is the same.  \Fm.

\emph{(e) Own-ack record.}  The Lean passes \nm{false} for \nm{acked}
(Clock:47) and keeps no record; the docstring of \nm{prepareClock} says
that the unique duty tick gives one ack per slot.  The specification keeps
$\AC.\acked$ across a restart (S21).  On a run with one tick at each time
the two agree.  \Fm{} under that condition, as ruled in S33.

\emph{(f) Output duty slot.}  \nm{outputHead} uses the current slot~$s$
as the duty slot of eligibility (Outputs:63--76).  S2 (line~476) uses
$s^*+1$.  Before $e_s$, $s^*=s-1$, and the two agree.  From $e_s$ on,
$s^*=s$, and S2 uses $s+1$: no accepted block has slot $s+1$ before
$t_{s+1}$, so in S2 the current-slot exception applies to no block, and
in the Lean it applies to the blocks of slot~$s$.  So the head of
\fname{ac\_head} and of \fname{ac\_to\_fg} can differ in
$[e_s,t_{s+1})$, for example at $a_r=d_{rR}$.  The specification follows
the Lean, as ruled in S35; the difference is against S2, not against the
Lean.

\emph{(g) Records.}  \nm{State.pcSlots} (State:47) counts the PC records;
\nm{decisionRecordSlots} and \nm{heldCompanionSlots} (Params:96--100)
give 13.  The specification lists the same 13 records.  \Fm.

\emph{(h) Atomic selection and own objects.}  \nm{selectDuties}
(Clock:58) and P \nm{act} select the four duties from one post-clock
state.  P \nm{recordOwn} (P:1188--1193) writes the send record and sends,
and the driver (Driver:574--580, Macro:194--197) runs those records and a
drain; the node admits its own duty object only at a delivery.  A PC that the
node assembles goes to \nm{on\_pc} at once (\nm{onAck}, Pools:64--67).
The specification now has the same rules (S45): its tick does not
process its own duty objects.  S2~\S1.1 processes them before the external receipts, so
just after an honest proposal tick the S2 head can be the new block while
the head of the Lean and of the specification is still the parent.

\emph{(i) The A2 batch.}  At the call (Clock:45--47), with $w=s-1$, A2
reads \nm{batch (s - 2)}: the votes of label $w-1$ received before
$d_{w-1}$ and resolved before $e_{w-1}$ (P:861--877); the batch is empty
for $s\leq1$, and \nm{ackEligible} uses no batch for label~0.  The
specification now reads the same view (S32): \fname{vote\_view} of slot
$w-1$ with the cutoffs $t_{w-1}+6\Delta$ and $t_{w-1}+2\Delta$.  The
docstring of \nm{ackEligible} (Gate:30--32) and S2 name the frozen and
carried batch (\nm{headBatch w}); the lead rules the docstring a defect.
A2 safety after a restart in $(t_{w-1},t_w]$ is OPEN.

\emph{(j) Slot~0.}  \nm{clockJoined} (Clock:10--12) is false at $s=0$, so
no honest node sends a slot-0 vote or runs a slot-0 duty.  The
specification now has the same test, and its $\AC.\restartTime$ starts at
$t_0=0$ (S34).  A faulty slot-0 vote is still admitted (\nm{onHead},
Pools:13--24), and a single early genesis vote can make a strict majority
of the participants of the raw test (Walks:43--48), so the raw result of
slot~0 can be non-$\bot$.  It has at most $a<\mu$ clean supporters, so FT
fails (Gate:17, 38--44): no honest node acks slot~0, and no valid slot-0
PC forms, although \nm{PC.valid} accepts label~0.  The difference is
against Goldilocks, which votes genesis at slot~0 (V24).

\emph{(k) Restart.}  The raw hook P \nm{restore} (P:1045--1051) applies
only \nm{restorePC}.  The run driver does the rest: \nm{macroRestore} and
\nm{restartWork} (Macro:171--186) run \nm{reconstruct} (the missed freezes
by \nm{lastFreeze}, from the sleep~$M$; Driver:383--398), a drain, the
replays of the non-genesis blocks of the active domain in depth order with
the restart stamp and the block relay (Driver:470--501, BlockCallback:
46--55), \nm{rebuild} (\nm{ac\_restart} with \nm{openLabels}: the scan of
every cumulative carrier into~$M$ and its open arm, the companion
normalization, and the regenerated requests; Driver:427--441,
Restart:75--83, Replay:21--68), and a last drain.  The specification (spec
3.9, S40) now has the same steps in the same order.  \Fm; the order of the
first drain against the replay relays is \UC.

\emph{(l) Opposition release.}  The specification and the Lean have the
same rule: \nm{authorityAncestor} gives the target when the slot of the
certificate, \nm{q.key.slot} (not the slot of the target block), is at
least the age slot, and \nm{latestAncestor} the deepest ancestor with
fewer than $\mu$ opponents otherwise (Authority:31--40).  D21c of
\nm{pcsa-decisions.md} has a different rule (R-OPP3); S39 lists the
differences.  The match of the specification with the Lean is checked;
the match with D21c fails.

\emph{(m) Dependency service.}  P keeps \nm{pending}, \nm{refused} and
\nm{origins} in \nm{Durable} (P:1018--1029).  P \nm{processReceipt}
(P:782--791) only calls the ordinary handler; the drain is in the run
driver (\nm{drainState}, Driver:363--370; \nm{drainOf} in
\nm{Macro.lean}:107--112), after each delivery, tick and restart step.  The
specification puts the same service in the PC-SA node (spec~4) and
proposes it as shared item (e).  The validation order: \nm{pcLeaf}
(Leaf:32--50) runs the stateless checks (with $B.\mathsf{parent}.\slot<B.\slot$),
the future test, the duplicate test, the parent test and the
parent-relative rows before the ancestry refusal, and the driver needs
\nm{rowChecks} before a leaf (Driver:219--225).  So a direct refusal never
marks an invalid block.  An inherited refusal (Driver:50--60, 79--85) is a
separate path and does not run the parent-relative row checks.  The
shared FG handler tests ancestry first; the PC-SA node of the
specification now has the Lean order for both paths (S46, V19).  With
that order the correspondence is \Fm; the exact interleaving of a drain
with the AC callbacks is \UC.

\emph{(n) Resolution.}  P \nm{resolveHeads} (P:620--625) tests only that
the head is held, and \nm{batchBefore}, \nm{batch} and \nm{releaseBatch}
(P:861--902) use that resolution; \nm{cleanCount} (Walks:23--27) reads the
resolved list as it is.  The specification now has no slot test in
resolution either (S41); its latest view and merged view keep the slot
test, as \nm{outputBatch} and \nm{headBatch}.  Against Goldilocks the
difference reaches only the unlocking: a block of slot $w+1$ is not
accepted before $t_{w+1}$ (Leaf:40), after the cut $e_w$ of \nm{batch},
but before the cut $q_w$ of \nm{releaseBatch}.  So a faulty label-$w$ vote
that names an accepted slot-$(w+1)$ block can be an opponent at the
unlocking.  \nm{oppositionCount} counts only clean seats, so an
equivocator is not an opponent.  The FG-SA draft of the main tree
(uncommitted) also counts the equivocators; this difference is recorded
in S41, to revisit when the FG-SA fast confirmation is committed.

\emph{(o) Population.}  S2~\S0 and the Lean (Params:14--15) have a PC
committee of $n$ seats and an FG population of $N$ validators, with
separate fault bounds; the claims (\nm{Regimes.lean}) use slot-dependent
holders.  The specification has one population, $n=N$, with every
committee~$V$, as \nm{shared/composition.tex}.  It is a specialization
(S42); the FG fault bound and the FG liveness quorum stay premises of the
transferred claims.

\emph{(p) Witness order.}  \nm{assemblePC} (Pools:29--37) orders the
acks by the decoded seat index of the slot.  The specification now does
the same (S43); this is a form difference against S6, which ordered the
acks by validator index.

\emph{Checked equal.}  These rules of the specification were compared
line by line with the Lean and have no difference: the order of
\nm{prepareClock} (carrier arm, unlock, update, early cut, freeze, ack
test, source); \nm{runGate} (raw result, FT, guard, record, lock writes
with the deadline $d_w+L$) apart from the path domain of V9;
\nm{ackEligible};
\nm{recentSource} and \nm{viewSource} with the label-0 empty freeze;
\nm{roleReceipt}, \nm{carrierReceipt} and \nm{acceptTargets};
\nm{historyReceipt}, \nm{addRecent}, \nm{recentCompatible} and
\nm{boundedHistoryCompatible}; \nm{targetStatus} and \nm{roleAuthority};
\nm{effectiveLock}; \nm{releaseLocks}; \nm{resolveHeads}; \nm{pcOutputs} (the clamp);
\nm{outputHead} at slot~0; \nm{onHead}, \nm{onAck}, and \nm{assemblePC};
\nm{clockJoined}; \nm{recordOwn};
\nm{on\_pc}; \nm{prepareSource} and \nm{ac\_to\_fg}; \nm{headDomain} and
\nm{headBatch}; P \nm{act} for the ack payload $(s-1,X.\mathsf{root},
X.\depth)$.

\section{Claims}\label{sec:claims}

The aggregate \nm{PCSAStage2Claims} (Claims:123) has twelve fields:
\nm{localModel}, \nm{healthySafety}, \nm{healthyInclusion},
\nm{formation}, \nm{recoveryService}, \nm{recoveryServiceValidator},
\nm{armedProtection}, \nm{localStrongProtection}, \nm{coveredSlowP1},
\nm{fgSafety}, \nm{fgFinality} and \nm{targetResolution}.  Its docstring
says: ``This checkpoint has no inhabitant.''  So no field is proved, and
each field is OPEN.  The claim table of spec~6 maps each S2 claim to a
field; the mapping is by name and by the S2 lines that the Lean
docstrings cite, and the rows marked \UC{} there are not checked.  The
fields range over canonical bindings of runs of \nm{pcProtocol}.  So a
proof of a field is about the rules of P.  Apart from loose BJ and the
weak E1, the specification differs from P only in items (b), (e) and (o), which are a
content abstraction, a form under one duty tick per time, and a
specialization.  Loose BJ, in the shared FG, changes the behaviour: a
justification at a bound height waits for the finality of its
predecessor, and the client reads its finality pair from the entry state
of its source (spec~6, paragraph Loose BJ).  P has neither, and no field
covers it; \nm{fgFinality} is a claim about the FG before loose BJ.
The weak E1 of the shared FG is not in P either: a timeout at~$h$ is E1
with a finality pair at~$h$ only when the two name different blocks, and
the client record keeps the set of the blocks that its timeouts name.
The Lean evidence predicate (\nm{conflictsWithFinality}) needs the root
test on its timeout arm.
The fields \nm{localModel.capacity} and
\nm{localModel.certificate\_provenance} (\nm{Local.lean}:67--105) state the
13-record and the provenance claims; \nm{formation} states the common key
(\nm{Formation.lean}:43--61).

The S2 claim table has REFUTED rows (spec~6).  No Lean field states a
REFUTED claim.

\section{Shared interface}\label{sec:shared}

Spec~5 lists the pending shared items: (a) the receipt frame, (b) the
atomic duty selection, (c) no pre-finality hook, (d) the historical read
service, (e) the dependency service, and (f) the relay.  The Lean has (a),
(b), (d), (e) and (f) inside the PC-SA run protocol P, its run driver
(\nm{Driver.lean}, \nm{Macro.lean}) and the shared \nm{FGCompose} layer
(\nm{RefusalEvidence}, \nm{ReceiptToken}).  S2
\S14.1 marks the shared refusal report as a REQUIRED rule that the shared
spec and the Lean implementation do not yet give.  This note makes no
change to the shared files.

\section{Legacy}\label{sec:legacy}

The earlier versions of this note (commits \nm{d8310c8}, \nm{fd996a5} and
\nm{f10fa3b}) compared the r5 text with the stage-1 model of the clone
\nm{~/lean/available-chain}.  That model (\nm{GridCO}, \nm{pcs},
\nm{pcStamps}, \nm{leaseRanks}, \nm{proposerPC}) is not the release
model, and those comparisons do not apply to the stage-2 specification.
The git history keeps them.

\end{document}
